\documentclass[11pt]{article}

\usepackage[margin=1in]{geometry}
\usepackage{amsmath,amssymb,amsthm,mathtools}
\usepackage{graphicx}
\usepackage[hidelinks]{hyperref}
\usepackage{enumitem}
\usepackage{microtype}

\setlength{\parindent}{0pt}
\setlength{\parskip}{0.55em}
\setlist[itemize]{leftmargin=1.45em,itemsep=0.2em}

\newtheorem{theorem}{Theorem}
\newtheorem{lemma}[theorem]{Lemma}
\newtheorem{corollary}[theorem]{Corollary}
\newtheorem{remark}[theorem]{Remark}

\newcommand{\C}{\mathbb C}
\newcommand{\R}{\mathbb R}
\newcommand{\N}{\mathbb N}
\newcommand{\cE}{\mathcal E}
\newcommand{\cF}{\mathcal F}
\newcommand{\dd}{\,d}
\newcommand{\pa}{\partial}
\newcommand{\ad}{A^{\dagger}}
\DeclareMathOperator{\im}{Im}

\begin{document}

\begin{center}
{\Large Full Solution Packet}\\[0.25em]
{\large A single line determines all twisted spherical means on $\C$}\\[0.45em]
\textbf{Status:} candidate full solution, likely valid\\
\textbf{Target:} the odd-Coxeter conjecture in arXiv:1103.4571
\end{center}

\section*{Source and exact conjecture}

The source is R. K. Srivastava, \emph{Coxeter system of lines are sets of
injectivity for the twisted spherical means on $\C$}, arXiv:1103.4571
\cite{SriLines}.  On page 14, after proving injectivity for Coxeter systems
with an even number of lines, the paper asks the following.

\begin{center}
\includegraphics[width=\textwidth]{figures/open_problem_crop.png}
\end{center}

In the source notation, if $\omega$ is an $N$th root of unity and
$\Sigma_N=\bigcup_{l=0}^{N-1}\omega^l\R$, the conjecture says that
for $f\in L^2(\C)$,
\[
 f\times\mu_r(\omega^l x)=0
 \quad(r>0,\ x\in\R,\ 0\le l<N)
 \qquad\Longrightarrow\qquad f=0\quad\text{a.e.}
\]
The surrounding discussion treats odd $N$ as the unresolved case.  The
result below proves the stronger assertion that the centers on just one line
through the origin suffice, and it works for every $L^p$ with $1\le p\le2$.

\section*{Definitions and normalization}

For $z,w\in\C$, write
\[
 (f\times g)(z)=\int_{\C}f(z-w)g(w)
 e^{\frac{i}{2}\im(z\bar w)}\dd w.
\]
If $\mu_r$ is normalized arc measure on $|w|=r$, then the twisted
spherical mean is
\[
 f\times\mu_r(z)=\int_{|w|=r}f(z-w)
 e^{\frac{i}{2}\im(z\bar w)}\dd\mu_r(w).
\]
Let
\[
 \varphi_k(z)=L_k\!\left(\frac{|z|^2}{2}\right)e^{-|z|^2/4},
 \qquad Q_k f=f\times\varphi_k.
\]
Harmless nonzero normalization constants in $\varphi_k$ do not affect any
vanishing statement.  The special Hermite operator is
\[
 \mathcal L=-\Delta+\frac{|z|^2}{4}
 -i\left(x\frac{\pa}{\pa y}-y\frac{\pa}{\pa x}\right),
 \qquad z=x+iy.
\]
The standard special Hermite expansion says that $Q_k f$ belongs to the
$\mathcal L$-eigenspace with eigenvalue $2k+1$, and the collection of these
projections reconstructs $f$ \cite{SriLines,GangulyThangavelu}.  We include
the part of this structure needed for the new uniqueness argument.

\section*{Result}

\begin{theorem}[Single-line injectivity]\label{thm:main}
Let $1\le p\le2$ and $f\in L^p(\C)$.  Suppose
\[
 f\times\mu_r(x)=0
 \qquad\text{for every }x\in\R\text{ and }r>0.
\]
Then $f=0$ almost everywhere.
\end{theorem}

\begin{corollary}[Odd Coxeter systems]\label{cor:coxeter}
Every line through the origin is a set of injectivity for twisted spherical
means on $L^p(\C)$, $1\le p\le2$.  Consequently every Coxeter system
$\Sigma_N$ is a set of injectivity, whether $N$ is odd or even.  In
particular, the conjecture on page 14 of arXiv:1103.4571 holds.
\end{corollary}

\section*{Proof intuition}

The spherical-mean hypothesis forces every spectral projection $Q_kf$ to
vanish on the real axis.  The missing point in the source is to show that
restriction to the axis is injective on each infinite-dimensional Landau
level.

Introduce annihilation and creation operators
\[
 A=\pa_z+\frac{\bar z}{4},
 \qquad
 \ad=-\pa_{\bar z}+\frac z4.
\]
Every vector in the $k$th level is $(\ad)^k$ applied to an
anti-holomorphic Fock vector $e^{-|z|^2/4}h(\bar z)$.  On the real axis,
the equation that this vector vanishes is a normal-ordered differential
equation
\[
 \sum_{j=0}^k\binom{k}{j}\left(\frac x2\right)^{k-j}
 (-1)^j h^{(j)}(x)=0.
\]
After $h(x)=e^{x^2/4}p(x)$, it becomes a constant-coefficient ODE
$P_k(D)p=0$, where
\[
 P_k(s)=2^{-k}i^kH_k(is).
\]
The zeros of the Hermite polynomial $H_k$ are real and simple, so $p$ is a
finite sum of exponentials with distinct real frequencies.  The Fock norm,
however, reduces to the $L^2$ norm of that exponential sum along every
horizontal line.  Such a sum has positive mean square unless it is zero.
Thus each $Q_kf$ is zero everywhere; special-Hermite reconstruction then
gives $f=0$.

\section*{The eigenspace restriction theorem}

Put
\[
 \cE_k=\ker_{L^2(\C)}(\mathcal L-(2k+1)I).
\]

\begin{lemma}[Creation representation]\label{lem:creation}
For every $v\in\cE_k$ there is an entire function $h$ such that
\begin{equation}\label{eq:creation}
 v(z)=c_k(\ad)^k\left(e^{-|z|^2/4}h(\bar z)\right),
 \qquad c_k\ne0,
\end{equation}
and
\begin{equation}\label{eq:fock}
 \int_{\C}|h(\bar z)|^2e^{-|z|^2/2}\dd z<\infty.
\end{equation}
\end{lemma}

\begin{proof}
A direct calculation gives
\[
 [A,\ad]=\frac12I,
 \qquad
 \mathcal L=4\ad A+I.
\]
Thus $A$ lowers the eigenvalue by $2$ and $\ad$ raises it by $2$.
If $v\in\cE_k$, then
\[
 \ad Av=\frac{k}{2}v.
\]
For $k>0$ this identity can be iterated, giving
\[
 v=\frac{2^k}{k!}(\ad)^kA^kv.
\]
Set $u=A^kv$.  Then $u\in\cE_0$, and the factorization gives
$4\|Au\|_2^2=\langle(\mathcal L-I)u,u\rangle=0$.  Thus $Au=0$, which is
\[
 \pa_z\left(e^{|z|^2/4}u(z)\right)=0.
\]
Hence $u(z)=e^{-|z|^2/4}h(\bar z)$ for an entire $h$.  Since $u\in L^2$,
its norm is exactly the integral in \eqref{eq:fock}.  This proves the
claim, including $k=0$.
\end{proof}

\begin{lemma}[A Landau level cannot vanish on a line]\label{lem:axis}
If $v\in\cE_k$ and $v(x)=0$ for every $x\in\R$, then $v=0$.
\end{lemma}

\begin{proof}
Use Lemma \ref{lem:creation}.  Writing $w=\bar z$, induction gives
\[
 (\ad)^k\left(e^{-|z|^2/4}h(w)\right)
 =e^{-|z|^2/4}\left(\frac z2-\pa_w\right)^kh(w),
\]
where $z$ is held fixed while $\pa_w$ acts.  Therefore restriction to
$z=w=x\in\R$ yields
\begin{equation}\label{eq:normal}
 T_kh(x):=\sum_{j=0}^k\binom{k}{j}
 \left(\frac x2\right)^{k-j}(-1)^jh^{(j)}(x)=0.
\end{equation}
The left side is entire, so the identity holds for complex $x$ as well.

Write $h(x)=e^{x^2/4}p(x)$.  The exponential generating function of the
operators in \eqref{eq:normal} is
\begin{align*}
 \sum_{k=0}^{\infty}\frac{t^k}{k!}T_kh(x)
 &=e^{tx/2}h(x-t)\\
 &=e^{x^2/4}e^{t^2/4}p(x-t).
\end{align*}
Consequently
\[
 T_kh(x)=e^{x^2/4}P_k(D)p(x),
 \quad D=\frac{d}{dx},
\]
where
\begin{equation}\label{eq:pk}
 P_k(s)=k!\sum_{r=0}^{\lfloor k/2\rfloor}
 \frac{(-s)^{k-2r}}{4^r r!(k-2r)!}
 =2^{-k}i^kH_k(is).
\end{equation}
For $k=0$, \eqref{eq:normal} immediately gives $h=0$.  Suppose $k\ge1$.
The physicists' Hermite polynomial $H_k$ has $k$ distinct real zeros
$\xi_1,\ldots,\xi_k$ \cite{Szego}.  Hence the roots of $P_k$ are the
distinct imaginary numbers $-i\xi_j$.  Every entire solution of
$P_k(D)p=0$ has the form
\begin{equation}\label{eq:exp-sum}
 p(w)=\sum_{j=1}^kc_je^{-i\xi_jw}.
\end{equation}

It remains to use the Fock condition, which is the decisive obstruction.
For $z=x+iy$ and $w=\bar z=x-iy$,
\begin{align*}
 e^{-|z|^2/2}|h(\bar z)|^2
 &=e^{-(x^2+y^2)/2}
   \left|e^{(x-iy)^2/4}p(x-iy)\right|^2\\
 &=e^{-y^2}|p(x-iy)|^2.
\end{align*}
For a fixed $y$, \eqref{eq:exp-sum} is a finite exponential sum in $x$:
\[
 p(x-iy)=\sum_{j=1}^k c_je^{-\xi_jy}e^{-i\xi_jx}.
\]
Distinct frequencies are orthogonal in mean, so
\[
 \lim_{R\to\infty}\frac1{2R}\int_{-R}^{R}|p(x-iy)|^2\dd x
 =\sum_{j=1}^k|c_j|^2e^{-2\xi_jy}.
\]
If one $c_j$ is nonzero, the right side is positive and the integral over
$x\in\R$ is infinite.  Tonelli's theorem then contradicts
\eqref{eq:fock}.  Thus all $c_j=0$, whence $h=0$ and $v=0$.
\end{proof}

\section*{Proof of the main theorem}

\begin{proof}[Proof of Theorem \ref{thm:main}]
Fix $k\ge0$.  Choose $s\in[1,2]$ from
\[
 1+\frac12=\frac1p+\frac1s.
\]
Because $\varphi_k$ is a polynomial times a Gaussian, it belongs to every
$L^s(\C)$.  The twisted Young inequality gives
\[
 \|Q_kf\|_2=\|f\times\varphi_k\|_2
 \le \|f\|_p\|\varphi_k\|_s<\infty.
\]
Thus $Q_kf\in\cE_k$.  Polar coordinates and Fubini give, with a nonzero
normalization constant $c$ independent of $x$,
\[
 Q_kf(x)=c\int_0^{\infty}
 (f\times\mu_r)(x)\varphi_k(r)r\dd r=0,
 \qquad x\in\R.
\]
The integral is absolutely legitimate: translation preserves $L^p$, while
the Gaussian factor belongs to the dual space.  Lemma \ref{lem:axis} now
implies
\[
 Q_kf=0\qquad\text{for every }k\ge0.
\]

For completeness, use the special-Hermite heat kernel.  The Laguerre
generating function gives
\[
 K_t(z)=\frac1{2\pi}\sum_{k=0}^{\infty}
 e^{-(2k+1)t}\varphi_k(z)
 =\frac{1}{4\pi\sinh t}
   \exp\!\left(-\frac14\coth(t)|z|^2\right).
\]
The series converges rapidly for $t>0$.  Hence, in tempered distributions
and in $L^p$,
\[
 e^{-t\mathcal L}f
 =f\times K_t
 =\frac1{2\pi}\sum_{k=0}^{\infty}e^{-(2k+1)t}Q_kf=0.
\]
Moreover, $K_t$ is a Gaussian concentrating at the origin and
$\int_{\C}K_t=1/\cosh t\to1$.  Strong continuity of twisted translations
therefore gives $f\times K_t\to f$ in $L^p(\C)$ as $t\downarrow0$ for every
$1\le p<\infty$ (the usual approximate-identity proof applies verbatim).
Therefore $f=0$.
\end{proof}

\begin{proof}[Proof of Corollary \ref{cor:coxeter}]
Twisted convolution with radial measures is covariant under rotations.
Applying Theorem \ref{thm:main} after a rotation proves injectivity for
every line $e^{i\theta}\R$.  Every Coxeter system $\Sigma_N$ contains at
least one such line, so vanishing for centers on $\Sigma_N$ implies vanishing
on that line and hence $f=0$.
\end{proof}

\section*{Verification and novelty check}

The exact checker
\texttt{code/check\_hermite\_axis\_ode.py} verifies for levels
$0\le k\le12$ that:
\begin{itemize}
\item coefficient extraction from $e^{t^2/4-ts}$ gives \eqref{eq:pk};
\item $P_k(s)=2^{-k}i^kH_k(is)$ exactly;
\item direct normal ordering on a generic polynomial agrees with
      $T_k(e^{x^2/4}p)=e^{x^2/4}P_k(D)p$; and
\item the computed Hermite zeros are real and distinct.
\end{itemize}
It was run with
\begin{center}
\texttt{conda run --no-capture-output -n sandbox python\\
code/check\_hermite\_axis\_ode.py --max-k 12}
\end{center}
and passed all levels.  This finite computation is not used as proof; the
generating-function and Hermite-zero argument above handles every $k$.

The bounded novelty search was performed on 26 August 2026.  The run's
registry, solution, attempt, and proof-gap indexes were searched for the
three relevant arXiv ids and for ``odd Coxeter'', ``twisted spherical means'',
and ``single line''; no prior packet was found.  The local parsed arXiv corpus
was searched for the title, arXiv:1103.4571, and the odd-line terminology.
The only directly relevant hits were the source, the adjacent papers
arXiv:1204.2773 and arXiv:1204.3076, and the later arXiv:1306.5658, which still
describes real-cone injectivity as unsolved.  Web searches restricted to
arXiv used the same phrases together with ``real axis'', ``Landau level'',
and ``polyanalytic Fock''.  They found background papers on polyanalytic Fock
spaces but no paper claiming the single-line theorem or the odd-Coxeter
conjecture.  This is a bounded search, not a proof of novelty.

\section*{Limitations and human-review recommendation}

The theorem is confined to $1\le p\le2$, exactly the range in which twisted
Young places every spectral projection in $L^2$.  It concerns lines through
the origin and does not settle the source's broader questions about arbitrary
real-analytic curves.  As usual, the hypothesis is interpreted wherever the
twisted spherical means are defined pointwise; the proof itself uses only the
resulting polar-integral identities.

The packet is recommended for expert review as a candidate full solution.
The highest-value checks are: (i) consistency of the twist convention with
$\mathcal L=4\ad A+I$ and the assignment of $Q_kf$ to $\cE_k$;
(ii) the normal-order generating function in Lemma \ref{lem:axis}; and
(iii) the $L^p$ special-Hermite heat-kernel reconstruction at $p=1$.

\begin{thebibliography}{9}

\bibitem{SriLines}
R. K. Srivastava,
\emph{Coxeter system of lines are sets of injectivity for the twisted
spherical means on $\C$}, arXiv:1103.4571.

\bibitem{SriExpansion}
R. K. Srivastava,
\emph{Real analytic expansion of spectral projection and extension of
Hecke--Bochner identity}, arXiv:1204.3076.

\bibitem{GangulyThangavelu}
P. Ganguly and S. Thangavelu,
\emph{Theorems of Chernoff and Ingham for certain eigenfunction expansions},
arXiv:2011.09940.

\bibitem{Szego}
G. Szeg\H{o}, \emph{Orthogonal Polynomials}, American Mathematical Society
Colloquium Publications, Vol. 23, 4th ed., 1975.

\end{thebibliography}

\end{document}
